% Generated draft: requires mathematical and source review.
\documentclass[11pt]{article}
\usepackage[margin=25mm]{geometry}
\usepackage{fontspec}
\setmainfont{Latin Modern Roman}
\usepackage{amsmath,amssymb,mathtools,unicode-math}
\setmathfont{Latin Modern Math}
\usepackage{enumitem,needspace,xurl,hyperref}
\setlength{\parindent}{0pt}
\setlength{\parskip}{5pt}
\setlength{\emergencystretch}{3em}
\newcommand{\N}{\mathbb{N}}
\newcommand{\Z}{\mathbb{Z}}
\newcommand{\Q}{\mathbb{Q}}
\newcommand{\R}{\mathbb{R}}
\newcommand{\C}{\mathbb{C}}
\newcommand{\F}{\mathbb{F}}
\newcommand{\A}{\mathbb{A}}
\newcommand{\PP}{\mathbb{P}}
\newcommand{\E}{\mathbb{E}}
\newcommand{\eps}{\varepsilon}
\newcommand{\dd}{\,\mathrm{d}}
\begin{document}
\section*{UnsolvedMath ID: 800981}
\textbf{Source problem code:} catalogue entry
\subsection{Definitions and mathematical statement}
Let \(T\) be the abstract \(\Theta\)-graph: two distinct vertices joined by three edges. The supplied question asks whether every metric on \(S^2\) admits a stationary realization of \(T\). It does not specify the class of metrics or the admissible realizations.

For a precise conditional formulation, suppose that \(g\) is a smooth Riemannian metric on \(S^2\). A nondegenerate regular realization consists of distinct points \(p,q\in S^2\) and three smooth regular curves \(\gamma_i:[0,1]\to S^2\), for \(i=1,2,3\), with \(\gamma_i(0)=p\) and \(\gamma_i(1)=q\). Its total length is
\[
L_g(\gamma_1,\gamma_2,\gamma_3)=\sum_{i=1}^{3}\int_0^1\lVert\dot\gamma_i(t)\rVert_g\,\dd t.
\]
Call this realization stationary if the first variation of \(L_g\) vanishes for every smooth variation through such realizations, allowing both common endpoints to move. Equivalently, each edge is a geodesic up to reparametrization, and the outgoing unit tangent vectors balance at each vertex:
\[
\sum_{i=1}^{3}\frac{\dot\gamma_i(0)}{\lVert\dot\gamma_i(0)\rVert_g}=0,
\qquad
\sum_{i=1}^{3}\frac{-\dot\gamma_i(1)}{\lVert\dot\gamma_i(1)\rVert_g}=0.
\]
Thus the three incident edges meet at pairwise angles of \(2\pi/3\) at each vertex.

An embedded realization additionally requires each \(\gamma_i\) to be injective and
\[
\gamma_i([0,1])\cap\gamma_j([0,1])=\{p,q\}\qquad(i\ne j).
\]
For either the regular-realization convention or the embedded-realization convention, the corresponding precise question is:
\[
\text{For every smooth Riemannian metric }g\text{ on }S^2,
\text{ does there exist a stationary realization of }T
\text{ satisfying that convention?}
\]
These are conditional formulations. The supplied record does not determine whether embeddedness is required, whether other degeneracies or intersections are allowed, or whether ``metric'' means a smooth Riemannian metric. Those qualifications must be resolved before selecting a definitive formulation.
\subsection{Short English statement}
Does every metric on the two-dimensional sphere admit a stationary \(\Theta\)-graph? Under a smooth Riemannian interpretation, stationarity means that the three edges are geodesic arcs whose outgoing unit tangents balance at both vertices. The supplied record leaves the metric class and the requirement of embeddedness unspecified.
\subsection{Sources}
\begin{itemize}
\item Local UnsolvedMath registry: permanent ID \(800981\), source problem code \texttt{catalogue entry}; supplied complete source record, upstream revision \texttt{d371ee917ac9647fbb06cd6fc529deec6f3dda82}.
\item Burns and Matveev, \emph{Open Problems and Questions about Geodesics}, Question 62, listed in the supplied record as 2019. Supplied source URL: \url{https://arxiv.org/abs/1308.5417}. The record reports access on 2026-10-16 and extraction by \texttt{arxiv\_latex}; these are supplied provenance details, not independently verified access or source assessment.
\item The title and background supply the attribution/context label \texttt{[Has-Mor1996]}. No full bibliographic entry or separate URL for that reference is supplied, so its authorship and precise relationship to the question remain unresolved.
\end{itemize}
\end{document}
